\chapter{实验补充材料}

\section{附录说明}
本附录给出正文之外的补充材料，包括关键超参数表、RCTA aggressive 配置摘要以及四个关键单源场景的完整结果。所有内容均服务于本文正文聚焦的单源到单目标跨工况诊断任务。

\section{基线方法关键超参数}
表~\ref{tab:appendix-baseline-hparams} 给出了本文 6 种基线方法在 aggressive 配置下采用的关键超参数。这些配置由 \texttt{configs/method/*.yaml} 与实验覆盖项共同确定。

\begin{table}[htbp]
\centering
\caption{基线方法关键超参数}
\label{tab:appendix-baseline-hparams}
\small
\begin{tabular}{lccc p{6.6cm}}
\hline
方法 & epochs & batch size & 学习率 & 关键对齐参数 \\
\hline
Source Only & 70 & 32 & $1\times 10^{-4}$ & 无显式对齐项 \\
CORAL & 70 & 32 & $1\times 10^{-4}$ & $\lambda_{coral}=0.5$，warm-start \\
DAN & 70 & 32 & $1\times 10^{-4}$ & $\lambda_{mmd}=0.05$，5 组高斯核尺度 \\
DANN & 70 & 32 & $1\times 10^{-4}$ & $\lambda_{adv}=0.5$，GRL warm-start，判别器 256$\times$2 \\
CDAN & 70 & 32 & $1\times 10^{-4}$ & $\lambda_{cadv}=0.2$，randomized dim=256，熵保护选择指标 \\
DeepJDOT & 70 & 32 & $1\times 10^{-4}$ & $\lambda_{ot}=1.0$，$\alpha=0.1$，$\beta=1.0$ \\
\hline
\end{tabular}
\normalsize
\end{table}

\section{RCTA aggressive 配置摘要}
表~\ref{tab:appendix-rcta-hparams} 汇总了本文单源四场景中采用的 RCTA aggressive 配置。该配置以 DANN 为基础对齐分支，同时融合 MCC 正则、可靠性门控、双原型记忆与一致性学习。

\begin{table}[htbp]
\centering
\caption{RCTA aggressive 配置摘要}
\label{tab:appendix-rcta-hparams}
\small
\begin{tabular}{l p{8.7cm}}
\hline
模块 & 配置 \\
\hline
训练轮数 & 96 \\
基础对齐分支 & DANN \\
MCC 正则 & $\lambda_{mcc}=0.1$，temperature=2.0 \\
教师网络 & EMA decay=0.99，teacher temperature=1.1 \\
可靠性权重 & cal\_conf=1.6，inv\_entropy=0.9，consistency=1.2，proto\_sim=1.4 \\
分数下界课程 & 0.24 $\rightarrow$ 0.58，持续 2400 steps \\
按类接受比例课程 & 0.30 $\rightarrow$ 0.92，持续 2400 steps \\
伪标签损失 & 最大权重 0.42，warmup 500 steps，可靠性加权，power=1.6 \\
原型损失 & 最大权重 0.08，start step=700，warmup 1000 steps \\
一致性损失 & 最大权重 0.12，start step=500，warmup 1000 steps \\
一致性门控 & 仅对 gate 内样本启用，reliability power=1.8 \\
数据增强 & 弱增强：轻微抖动/缩放；强增强：抖动、缩放、时间屏蔽、通道 dropout \\
\hline
\end{tabular}
\normalsize
\end{table}

\section{四个关键单源场景的完整结果}
表~\ref{tab:appendix-single-source-full} 给出了本文四个关键单源场景的完整结果，包括源域训练准确率、源域评估准确率、目标域准确率、目标域平衡准确率和最终选择 epoch。该表可作为正文统计结果的完整支撑材料。

\begin{table}[htbp]
\centering
\caption{四个关键单源场景的完整结果}
\label{tab:appendix-single-source-full}
\scriptsize
\begin{tabular}{lcccccc}
\hline
方法 & 场景 & 源域训练准确率 & 源域评估准确率 & 目标域准确率 & 目标域平衡准确率 & 选择 epoch \\
\hline
Source Only & mode2$\rightarrow$mode5 & 0.8845 & 0.8889 & 0.6094 & 0.6116 & 70 \\
CORAL & mode2$\rightarrow$mode5 & 0.8661 & 0.8571 & 0.6875 & 0.6871 & 69 \\
DAN & mode2$\rightarrow$mode5 & 0.8494 & 0.8536 & 0.8177 & 0.8190 & 68 \\
DANN & mode2$\rightarrow$mode5 & 0.8525 & 0.8519 & 0.8299 & 0.8310 & 69 \\
CDAN & mode2$\rightarrow$mode5 & 0.8661 & 0.8713 & 0.8177 & 0.8190 & 70 \\
DeepJDOT & mode2$\rightarrow$mode5 & 0.7853 & 0.7813 & 0.7135 & 0.7155 & 65 \\
RCTA & mode2$\rightarrow$mode5 & 0.8990 & 0.9048 & 0.8542 & 0.8552 & 94 \\
\hline
Source Only & mode3$\rightarrow$mode6 & 0.8759 & 0.8774 & 0.7064 & 0.7069 & 68 \\
CORAL & mode3$\rightarrow$mode6 & 0.8849 & 0.8774 & 0.7962 & 0.7966 & 65 \\
DAN & mode3$\rightarrow$mode6 & 0.8853 & 0.8843 & 0.8187 & 0.8190 & 59 \\
DANN & mode3$\rightarrow$mode6 & 0.8772 & 0.8705 & 0.8566 & 0.8569 & 69 \\
CDAN & mode3$\rightarrow$mode6 & 0.8681 & 0.8756 & 0.8204 & 0.8207 & 62 \\
DeepJDOT & mode3$\rightarrow$mode6 & 0.8203 & 0.8117 & 0.7168 & 0.7172 & 70 \\
RCTA & mode3$\rightarrow$mode6 & 0.9060 & 0.9102 & 0.8756 & 0.8759 & 96 \\
\hline
Source Only & mode5$\rightarrow$mode2 & 0.8691 & 0.8767 & 0.5203 & 0.5190 & 69 \\
CORAL & mode5$\rightarrow$mode2 & 0.8427 & 0.8403 & 0.6808 & 0.6862 & 66 \\
DAN & mode5$\rightarrow$mode2 & 0.8704 & 0.8663 & 0.8131 & 0.8052 & 67 \\
DANN & mode5$\rightarrow$mode2 & 0.8435 & 0.8472 & 0.7831 & 0.7766 & 66 \\
CDAN & mode5$\rightarrow$mode2 & 0.8721 & 0.8750 & 0.7478 & 0.7414 & 68 \\
DeepJDOT & mode5$\rightarrow$mode2 & 0.7811 & 0.7812 & 0.7072 & 0.7091 & 70 \\
RCTA & mode5$\rightarrow$mode2 & 0.9142 & 0.8872 & 0.8183 & 0.8184 & 96 \\
\hline
Source Only & mode6$\rightarrow$mode3 & 0.8654 & 0.8687 & 0.8463 & 0.8465 & 69 \\
CORAL & mode6$\rightarrow$mode3 & 0.8624 & 0.8618 & 0.8359 & 0.8362 & 68 \\
DAN & mode6$\rightarrow$mode3 & 0.8382 & 0.8342 & 0.7979 & 0.7983 & 70 \\
DANN & mode6$\rightarrow$mode3 & 0.8619 & 0.8618 & 0.8618 & 0.8621 & 70 \\
CDAN & mode6$\rightarrow$mode3 & 0.7934 & 0.7927 & 0.7547 & 0.7550 & 36 \\
DeepJDOT & mode6$\rightarrow$mode3 & 0.7994 & 0.7720 & 0.7478 & 0.7481 & 70 \\
RCTA & mode6$\rightarrow$mode3 & 0.9025 & 0.8981 & 0.8877 & 0.8879 & 95 \\
\hline
\end{tabular}
\normalsize
\end{table}

\section{补充说明}
正文中的汇总图由四个关键单源场景的比较结果自动生成，代表性定性图则选取自 mode2$\rightarrow$mode5 场景。之所以选择该场景，一方面是因为其相对 Source Only 的改进幅度较大，另一方面也因为该场景中 Source Only 与 RCTA 的结构差异更容易从混淆矩阵和 t-SNE 图中观察到。
